\documentclass[11pt,a4paper]{ctexart}

\usepackage[margin=1in]{geometry}
\usepackage{amsmath,amssymb,amsthm,mathtools,bm}
\usepackage{algorithm}
\usepackage{algpseudocode}
\usepackage{hyperref}
\usepackage{enumitem}
\usepackage{booktabs}
\usepackage{array}

\hypersetup{
  colorlinks=true,
  linkcolor=blue,
  urlcolor=blue,
  citecolor=blue
}

% --------------------------------------------------
% Macros
% --------------------------------------------------
\newcommand{\RR}{\mathbb{R}}
\newcommand{\ZZ}{\mathbb{Z}}
\newcommand{\CC}{\mathbb{C}}

\newcommand{\Rbig}{\mathcal{R}_N^{\RR}}
\newcommand{\Rint}{\mathcal{R}_N^{\ZZ}}
\newcommand{\Ssmall}{\mathcal{S}_n}

\newcommand{\Enc}{\mathsf{Enc}}
\newcommand{\Auto}[1]{\mathsf{Auto}_{#1}}
\newcommand{\Add}{\mathsf{Add}}
\newcommand{\Sub}{\mathsf{Sub}}
\newcommand{\MulPt}[1]{\mathsf{MulPt}_{#1}}

\newcommand{\Round}{\operatorname{Round}}
\newcommand{\rep}{\mathrel{\vDash_{\Delta}}}

\newcommand{\ctP}{\mathbf{ct}_{P}}
\newcommand{\ctm}{\mathbf{ct}_{-}}
\newcommand{\cte}{\mathbf{ct}_{\mathrm{even}}^{\uparrow}}
\newcommand{\ctoy}{\mathbf{ct}_{\mathrm{odd},Y}^{\uparrow}}
\newcommand{\cto}{\mathbf{ct}_{\mathrm{odd}}^{\uparrow}}

\newcommand{\Pe}{P_{\mathrm{even}}}
\newcommand{\Po}{P_{\mathrm{odd}}}

% --------------------------------------------------
% Theorem styles
% --------------------------------------------------
\newtheorem{definition}{定义}[section]
\newtheorem{theorem}[definition]{定理}
\newtheorem{remark}[definition]{注}
\newtheorem{proposition}[definition]{命题}

% --------------------------------------------------
% Algorithmic style tweaks
% --------------------------------------------------
\algrenewcommand\algorithmicrequire{\textbf{输入:}}
\algrenewcommand\algorithmicensure{\textbf{输出:}}
\algrenewcommand\algorithmiccomment[1]{\hfill$\triangleright$~#1}

% --------------------------------------------------
% Title
% --------------------------------------------------
\title{在 CKKS 密文上实现多项式的 FFT 式奇偶拆分与重构\\
\large 第二版：问题、主结果、正确性定理与伪代码}
\author{}
\date{}

\begin{document}


\maketitle


\tableofcontents
\newpage

\begin{abstract}
本文给出一份面向实现的第二版整理。我们固定
\[
P(Y)=\sum_{j=0}^{2n-1}p_jY^j \in \RR[Y]/(Y^{2n}+1),
\]
并研究如何在 CKKS 密文上直接完成 FFT 式奇偶拆分：
\[
P(Y)=P_{\mathrm{even}}(Y^2)+Y\,P_{\mathrm{odd}}(Y^2),
\]
以及对应的反向重构。与第一版不同，本文不再重复代数证明，而是采用
\[
\text{Problem} \;\to\; \text{Main Result} \;\to\; \text{Correctness Theorem} \;\to\; \text{Algorithm}
\]
的组织方式，给出可直接转译到程序中的伪代码。文末附上 OpenFHE、Microsoft SEAL 与 Lattigo v5 的 API 对应说明，并指出实现时必须特别注意的自同构索引、plaintext 构造、level/scale 对齐等事项。
\end{abstract}

\section{Problem}

\subsection{目标}

固定一个 2 的幂 \(N\ge 1\)，以及一个满足 \(n\mid N\) 且 \(n<N\) 的 2 的幂。令
\[
k=\frac{N}{2n}\in \ZZ_{>0},
\qquad
Y=X^k.
\]
于是
\[
Y^{2n}=X^{2nk}=X^N=-1.
\]

给定多项式
\[
P(Y)=\sum_{j=0}^{2n-1}p_jY^j\in \RR[Y]/(Y^{2n}+1),
\]
定义其 FFT 式奇偶分支为
\begin{equation}
\Pe(T)=\sum_{i=0}^{n-1}p_{2i}T^i,
\qquad
\Po(T)=\sum_{i=0}^{n-1}p_{2i+1}T^i.
\label{eq:def-even-odd}
\end{equation}

本文考虑如下两个问题：

\begin{enumerate}[label=\textbf{P\arabic*.}]
    \item \textbf{正向拆分问题：} 给定一个 CKKS 密文 \(\ctP\)，其语义上表示
    \[
    P(X^k)\in \Rbig=\RR[X]/(X^N+1),
    \]
    如何同态地产生两个分支密文，分别表示
    \[
    \Pe(X^{2k})
    \quad\text{和}\quad
    \Po(X^{2k})
    \]
    （或者表示 \(X^k\Po(X^{2k})\) 的非归一化奇分支）？

    \item \textbf{反向重构问题：} 给定上述两个分支密文，如何同态地恢复出表示 \(P(X^k)\) 的原始密文？
\end{enumerate}

\subsection{为什么选多项式域路线而不是系数向量矩阵路线}

虽然在抽象线性代数层面，可以把
\[
(p_0,p_1,\dots,p_{2n-1})^\top
\]
视为一个长度 \(2n\) 的系数向量，并用抽取矩阵 \(E,O\) 去得到偶、奇系数子向量；但在标准 CKKS 的高层接口中，最自然的对象是：

\begin{itemize}
    \item \textbf{槽向量表示}（经编码后的复向量）；
    \item \textbf{多项式明文表示}（环中的多项式）。
\end{itemize}

如果走 \(E,O\) 的路线，通常需要先做一次 CoeffToSlot，把系数信息显式搬到槽里；这会引入额外的线性变换成本。本文因此采用更直接的多项式域恒等式：
\begin{equation}
P(Y)=\Pe(Y^2)+Y\,\Po(Y^2), \label{eq:split}
\end{equation}
\begin{equation}
P(-Y)=\Pe(Y^2)-Y\,\Po(Y^2). \label{eq:split-minus}
\end{equation}

\section{Notation and CKKS Semantics}

\subsection{环与嵌入}

记
\[
\Ssmall=\RR[Y]/(Y^{2n}+1),
\qquad
\Rbig=\RR[X]/(X^N+1),
\qquad
\Rint=\ZZ[X]/(X^N+1).
\]
通过 \(Y=X^k\) 把 \(\Ssmall\) 嵌入 \(\Rbig\)。

\subsection{密文语义记号}

我们使用如下标准化记号来表达 CKKS 密文在 scale \(\Delta\) 下所表示的消息。

\begin{definition}[消息表示关系]
若一个 CKKS 密文 \(\mathbf{ct}\) 解密后对应的明文多项式可写成
\[
m(X)=\Round\!\bigl(\Delta F(X)\bigr)+e(X)\in \Rint,
\]
其中 \(F(X)\in \Rbig\)，而 \(e(X)\) 是通常的小噪声项，则记为
\[
\mathbf{ct}\rep F(X).
\]
\end{definition}

该记号只关注“被表示的实系数消息” \(F(X)\)，不显式追踪噪声、RNS 分解、模链变化、rescale 细节。若实现中需要进行 level/scale 对齐，则默认按标准 CKKS 方式处理。

\subsection{本文使用的同态算子}

本文只用到以下四类标准算子：

\begin{enumerate}[label=(\roman*)]
    \item \textbf{自同构：}
    \[
    \Auto{a}(\mathbf{ct}),
    \qquad a \text{ 为奇数},
    \]
    对应
    \[
    \sigma_a:\Rbig\to\Rbig,\qquad f(X)\mapsto f(X^a).
    \]

    \item \textbf{加法：}
    \[
    \Add(\mathbf{ct}_1,\mathbf{ct}_2).
    \]

    \item \textbf{减法：}
    \[
    \Sub(\mathbf{ct}_1,\mathbf{ct}_2).
    \]

    \item \textbf{明文--密文乘法：}
    \[
    \MulPt{u(X)}(\mathbf{ct}),
    \qquad u(X)\in \Rbig.
    \]
\end{enumerate}

\subsection{上标 \(\uparrow\) 的含义}

本文中的
\[
\cte,\qquad \ctoy,\qquad \cto
\]
都带有上标 \({}^{\uparrow}\)。这表示：对应消息虽然仍然存放在大环 \(\Rbig\) 中，但只含 \(X^{2k}\) 的幂，因此实际上来自一个更小子环的“提升表示”。

\section{Main Result}

\subsection{核心恒等式}

由 \eqref{eq:split} 和 \eqref{eq:split-minus} 立刻得到
\begin{equation}
\Pe(Y^2)=\frac{P(Y)+P(-Y)}{2}, \label{eq:even-formula}
\end{equation}
\begin{equation}
Y\,\Po(Y^2)=\frac{P(Y)-P(-Y)}{2}. \label{eq:odd-formula-y}
\end{equation}
若进一步右乘 \(Y^{-1}\)，则
\begin{equation}
\Po(Y^2)=Y^{-1}\cdot \frac{P(Y)-P(-Y)}{2}. \label{eq:odd-formula}
\end{equation}

在 \(Y=X^k\) 的嵌入下，上式变成
\begin{equation}
\Pe(X^{2k})=\frac{P(X^k)+P(-X^k)}{2}, \label{eq:even-formula-X}
\end{equation}
\begin{equation}
X^k\Po(X^{2k})=\frac{P(X^k)-P(-X^k)}{2}, \label{eq:odd-formula-X-y}
\end{equation}
\begin{equation}
\Po(X^{2k})=(X^k)^{-1}\cdot \frac{P(X^k)-P(-X^k)}{2}. \label{eq:odd-formula-X}
\end{equation}

\subsection{实现 \(Y\mapsto -Y\) 的特定自同构}

取
\begin{equation}
a=1+2n. \label{eq:a-index}
\end{equation}
则 \(a\) 为奇数，并且对应自同构 \(\sigma_a(X)=X^a\)。它满足
\begin{equation}
\sigma_a(X^k)=X^{k(1+2n)}=X^{k+N}=X^kX^N=-X^k. \label{eq:auto-y}
\end{equation}

因此，若 \(\ctP\rep P(X^k)\)，则
\[
\ctm:=\Auto{1+2n}(\ctP)
\]
满足
\[
\ctm \rep P(-X^k).
\]

\section{Correctness Theorems (No Proofs)}

本节只陈述结论，不再给证明。

\begin{theorem}[正向拆分：非归一化奇分支]
设
\[
\ctP\rep P(X^k).
\]
定义
\[
\ctm:=\Auto{1+2n}(\ctP),
\]
\[
\cte:=\MulPt{\frac12}\!\bigl(\Add(\ctP,\ctm)\bigr),
\]
\[
\ctoy:=\MulPt{\frac12}\!\bigl(\Sub(\ctP,\ctm)\bigr).
\]
则
\[
\cte \rep \Pe(X^{2k}),
\qquad
\ctoy \rep X^k\Po(X^{2k}).
\]
\end{theorem}

\begin{theorem}[正向拆分：归一化奇分支]
在上一结论的基础上，再定义
\[
\cto:=\MulPt{-X^{N-k}}(\ctoy).
\]
则
\[
\cto \rep \Po(X^{2k}).
\]
\end{theorem}

\begin{theorem}[反向重构：从非归一化奇分支]
若
\[
\cte \rep \Pe(X^{2k}),
\qquad
\ctoy \rep X^k\Po(X^{2k}),
\]
则
\[
\Add(\cte,\ctoy)\rep P(X^k).
\]
\end{theorem}

\begin{theorem}[反向重构：从归一化奇分支]
若
\[
\cte \rep \Pe(X^{2k}),
\qquad
\cto \rep \Po(X^{2k}),
\]
则
\[
\Add\!\bigl(\cte,\MulPt{X^k}(\cto)\bigr)\rep P(X^k).
\]
\end{theorem}

\section{Algorithms}

\subsection{约定}

下列伪代码默认：

\begin{itemize}
    \item 所有密文在进入某一步前都已完成必要的 \texttt{level/scale} 对齐；
    \item \texttt{Auto}\((1+2n)\) 所需的自同构评估密钥已预生成；
    \item \(\frac12\)、\(X^k\)、\(-X^{N-k}\) 这些明文乘子都能被构造为可参与 plaintext--ciphertext 乘法的明文对象。
\end{itemize}

\subsection{正向拆分：保留奇分支前面的 \(X^k\)}

\begin{algorithm}[H]
\caption{ForwardSplitKeepY}
\begin{algorithmic}[1]
\Require
A ciphertext \(\ctP\) such that \(\ctP \rep P(X^k)\), where
\[
P(Y)=\sum_{j=0}^{2n-1}p_jY^j,\qquad k=\frac{N}{2n}.
\]
\Ensure
Two ciphertexts \(\cte,\ctoy\) such that
\[
\cte \rep \Pe(X^{2k}),\qquad \ctoy \rep X^k\Po(X^{2k}).
\]
\State \(a \gets 1+2n\)
\State \(\ctm \gets \Auto{a}(\ctP)\)
\Comment{\(\ctm \rep P(-X^k)\)}
\State \(\mathbf{s}_{+} \gets \Add(\ctP,\ctm)\)
\State \(\mathbf{s}_{-} \gets \Sub(\ctP,\ctm)\)
\State \(\cte \gets \MulPt{\frac12}(\mathbf{s}_{+})\)
\Comment{\(\cte \rep \Pe(X^{2k})\)}
\State \(\ctoy \gets \MulPt{\frac12}(\mathbf{s}_{-})\)
\Comment{\(\ctoy \rep X^k\Po(X^{2k})\)}
\State \Return \((\cte,\ctoy)\)
\end{algorithmic}
\end{algorithm}

\subsection{正向拆分：去掉奇分支前面的 \(X^k\)}

\begin{algorithm}[H]
\caption{ForwardSplitNormalized}
\begin{algorithmic}[1]
\Require
A ciphertext \(\ctP\) such that \(\ctP \rep P(X^k)\).
\Ensure
Two ciphertexts \(\cte,\cto\) such that
\[
\cte \rep \Pe(X^{2k}),\qquad \cto \rep \Po(X^{2k}).
\]
\State \((\cte,\ctoy) \gets \textsc{ForwardSplitKeepY}(\ctP)\)
\State \(u(X)\gets -X^{N-k}\)
\Comment{Since \((X^k)^{-1}=-X^{N-k}\)}
\State \(\cto \gets \MulPt{u(X)}(\ctoy)\)
\Comment{\(\cto \rep \Po(X^{2k})\)}
\State \Return \((\cte,\cto)\)
\end{algorithmic}
\end{algorithm}

\subsection{反向重构：输入保留 \(X^k\) 的奇分支}

\begin{algorithm}[H]
\caption{MergeFromKeepY}
\begin{algorithmic}[1]
\Require
Two ciphertexts \(\cte,\ctoy\) such that
\[
\cte \rep \Pe(X^{2k}),\qquad \ctoy \rep X^k\Po(X^{2k}).
\]
\Ensure
A ciphertext \(\ctP\) such that
\[
\ctP \rep P(X^k).
\]
\State \(\ctP \gets \Add(\cte,\ctoy)\)
\State \Return \(\ctP\)
\end{algorithmic}
\end{algorithm}

\subsection{反向重构：输入为归一化奇分支}

\begin{algorithm}[H]
\caption{MergeFromNormalized}
\begin{algorithmic}[1]
\Require
Two ciphertexts \(\cte,\cto\) such that
\[
\cte \rep \Pe(X^{2k}),\qquad \cto \rep \Po(X^{2k}).
\]
\Ensure
A ciphertext \(\ctP\) such that
\[
\ctP \rep P(X^k).
\]
\State \(\mathbf{t}\gets \MulPt{X^k}(\cto)\)
\Comment{\(\mathbf{t}\rep X^k\Po(X^{2k})\)}
\State \(\ctP \gets \Add(\cte,\mathbf{t})\)
\State \Return \(\ctP\)
\end{algorithmic}
\end{algorithm}

\subsection{一个最小执行模板}

如果实现中需要明确地把“生成自同构密钥”也写入流程，那么最小模板可写成：

\begin{algorithm}[H]
\caption{EndToEndSplitAndMerge}
\begin{algorithmic}[1]
\Require
A ciphertext \(\ctP \rep P(X^k)\); secret-dependent evaluation-key generation is available offline.
\Ensure
A reconstructed ciphertext \(\widehat{\ctP}\rep P(X^k)\).
\State \(a\gets 1+2n\)
\State Generate / load the automorphism key for \(\Auto{a}\)
\State \((\cte,\cto)\gets \textsc{ForwardSplitNormalized}(\ctP)\)
\State \(\widehat{\ctP}\gets \textsc{MergeFromNormalized}(\cte,\cto)\)
\State \Return \(\widehat{\ctP}\)
\end{algorithmic}
\end{algorithm}

\section{Implementation Notes}

\subsection{总则}

虽然本文给出的算法在代数层面非常简洁，但把它落到具体库时，至少要注意四件事：

\begin{enumerate}[label=(\arabic*)]
    \item \textbf{自同构索引的表示方式：}  
    数学上我们写的是 \(\sigma_{1+2n}\)，即
    \[
    X\mapsto X^{1+2n}.
    \]
    但不同库对 Galois 元素、rotation index、automorphism index 的表示方式不同。实现时必须确认：你传入的“索引 / Galois element”确实实现了
    \[
    X^k \longmapsto -X^k.
    \]

    \item \textbf{plaintext 的构造方式：}  
    本文用到的明文乘子是
    \[
    \frac12,\qquad X^k,\qquad -X^{N-k}.
    \]
    常数 \(1/2\) 通常最容易；而单项式 \(X^k\)、\(-X^{N-k}\) 往往需要底层的 plaintext-polynomial 构造路径。如果库的高层 CKKS 接口只暴露“槽编码器”，那么实现这一步可能需要转到更底层的 plaintext 对象或自定义构造函数。

    \item \textbf{level/scale 对齐：}  
    做
    \[
    \ctP\pm \ctm
    \]
    之前，必须保证二者在同一 level、同一 scale（或者能通过库的自动机制对齐）。做 plaintext multiplication 之后，是否需要 rescale，取决于库的语义和你给 plaintext 指定的 scale。

    \item \textbf{系数域语义与槽语义：}  
    本文算法在“多项式消息语义”下描述。如果库层面默认以槽语义暴露 CKKS plaintext/ciphertext，那么你需要清楚自己是在使用高层编码接口，还是在直接操纵环里的 plaintext polynomial。
\end{enumerate}

\subsection{OpenFHE}

在 OpenFHE 中，可将本文四类抽象操作大致映射为：

\begin{center}
\begin{tabular}{>{\raggedright\arraybackslash}p{0.28\textwidth} >{\raggedright\arraybackslash}p{0.62\textwidth}}
\toprule
抽象操作 & OpenFHE 中的对应接口 \\ \midrule
生成 automorphism key &
\texttt{EvalAutomorphismKeyGen(privateKey, indexList)} \\[4pt]
应用 automorphism &
\texttt{EvalAutomorphism(ciphertext, i, evalKeyMap)} \\[4pt]
密文加法 / 减法 &
\texttt{EvalAdd}, \texttt{EvalSub} \\[4pt]
明文--密文乘法 &
\texttt{EvalMult(ciphertext, scalar)} 或 \texttt{EvalMult(ciphertext, plaintext)} \\[4pt]
\bottomrule
\end{tabular}
\end{center}

\paragraph{建议写法}
若你已经知道目标 automorphism index 对应数学上的 \(\sigma_{1+2n}\)，那么正向拆分最接近的实现骨架是：
\begin{align*}
\ctm &\leftarrow \texttt{cc->EvalAutomorphism}(\ctP,\; a,\; \texttt{evalAutoKeys}),\\
\cte &\leftarrow \texttt{cc->EvalMult}\bigl(\texttt{cc->EvalAdd}(\ctP,\ctm),\; 0.5\bigr),\\
\ctoy &\leftarrow \texttt{cc->EvalMult}\bigl(\texttt{cc->EvalSub}(\ctP,\ctm),\; 0.5\bigr).
\end{align*}

\paragraph{注意}
OpenFHE 的高层 CKKS 工作流通常以 packed plaintext 为中心。如果你要真正实现
\[
\MulPt{X^k}(\cdot),\qquad \MulPt{-X^{N-k}}(\cdot),
\]
需要确认你所用版本是否方便构造“系数域明文多项式”作为 \texttt{Plaintext}；否则更稳妥的办法是自己在较低层构造该明文对象，再交给 \texttt{EvalMult}。

\subsection{Microsoft SEAL}

在 SEAL 中，本文抽象操作可大致映射为：

\begin{center}
\begin{tabular}{>{\raggedright\arraybackslash}p{0.28\textwidth} >{\raggedright\arraybackslash}p{0.62\textwidth}}
\toprule
抽象操作 & Microsoft SEAL 中的对应接口 \\ \midrule
生成 Galois keys &
\texttt{KeyGenerator::create\_galois\_keys(...)} \\[4pt]
应用 automorphism &
\texttt{Evaluator::apply\_galois\_inplace(...)} \\[4pt]
密文加法 / 减法 &
\texttt{Evaluator::add\_inplace}, \texttt{Evaluator::sub\_inplace} \\[4pt]
明文--密文乘法 &
\texttt{Evaluator::multiply\_plain\_inplace(...)} \\[4pt]
plaintext 预处理 &
\texttt{Evaluator::transform\_to\_ntt\_inplace(...)}（若需要） \\[4pt]
\bottomrule
\end{tabular}
\end{center}

\paragraph{建议写法}
SEAL 中通常会把 \(\ctP\) 复制到临时对象，再原地修改。例如：
\begin{align*}
\ctm &\leftarrow \ctP;\\
&\texttt{evaluator.apply\_galois\_inplace}(\ctm,\; \texttt{galois\_elt},\; \texttt{galois\_keys});\\
\cte &\leftarrow \ctP;\\
&\texttt{evaluator.add\_inplace}(\cte,\ctm);\\
&\texttt{evaluator.multiply\_plain\_inplace}(\cte,\texttt{plain\_half});\\
\ctoy &\leftarrow \ctP;\\
&\texttt{evaluator.sub\_inplace}(\ctoy,\ctm);\\
&\texttt{evaluator.multiply\_plain\_inplace}(\ctoy,\texttt{plain\_half}).
\end{align*}

\paragraph{注意}
SEAL 对 CKKS plaintext 的 NTT 形式、\texttt{parms\_id}、scale 对齐非常严格。  
尤其是在做
\[
\texttt{multiply\_plain\_inplace}
\]
前，往往需要确保 plaintext 已经处于正确的 level，并在必要时通过
\[
\texttt{mod\_switch\_to\_inplace},\qquad \texttt{transform\_to\_ntt\_inplace}
\]
进行预处理。若要实现单项式 \(X^k\) 或 \(-X^{N-k}\) 的明文乘法，通常也需要手工构造对应的 \texttt{Plaintext}。

\subsection{Lattigo v5}

Lattigo v5 把 automorphism 更明确地暴露在 RLWE 层，而 CKKS 层的 evaluator 负责常规算术。可以大致映射为：

\begin{center}
\begin{tabular}{>{\raggedright\arraybackslash}p{0.28\textwidth} >{\raggedright\arraybackslash}p{0.62\textwidth}}
\toprule
抽象操作 & Lattigo v5 中的对应接口 \\ \midrule
生成 Galois key &
\texttt{rlwe.KeyGenerator.GenGaloisKey} / \texttt{GenGaloisKeyNew} \\[4pt]
应用 automorphism &
\texttt{rlwe.Evaluator.Automorphism} \\[4pt]
获取某些标准 Galois 元素 &
\texttt{ckks.Parameters.GaloisElementForRotation(...)} \\[4pt]
密文加法 / 减法 &
\texttt{ckks.Evaluator.Add}, \texttt{ckks.Evaluator.Sub} \\[4pt]
明文--密文乘法 &
\texttt{ckks.Evaluator.Mul} \\[4pt]
rescale &
\texttt{ckks.Evaluator.Rescale} \\[4pt]
\bottomrule
\end{tabular}
\end{center}

\paragraph{建议写法}
如果你已知目标 Galois element \(\texttt{galEl}\) 实现了
\[
X^k\mapsto -X^k,
\]
那么可在 RLWE 层先做 automorphism，再回到 CKKS evaluator 做算术。一个自然的分层写法是：

\begin{enumerate}[label=(\arabic*)]
    \item 用 \texttt{rlwe.KeyGenerator.GenGaloisKeyNew(galEl, sk, ...)} 生成所需 Galois key；
    \item 用 \texttt{rlwe.Evaluator.Automorphism(ctIn, galEl, ctOut)} 得到 \(\ctm\)；
    \item 用 \texttt{ckks.Evaluator.Add/Sub/Mul} 完成
    \[
    \cte,\;\ctoy,\;\cto
    \]
    的构造；
    \item 若 plaintext multiplication 改变了 scale，则按需要调用 \texttt{Rescale}。
\end{enumerate}

\paragraph{注意}
Lattigo 的 \texttt{GaloisElementForRotation(k)} 是“标准槽旋转”的 helper；而本文需要的是“在多项式环中实现 \(X^k\mapsto -X^k\) 的 automorphism”。这两者不应混淆。若你不能把目标 automorphism 等价地转化为一个标准 rotation，那么更稳妥的做法是直接在 RLWE 层使用所需的 \texttt{galEl}。

\section{Practical Checklist}

下面给出一个真正实现时的检查清单。

\begin{enumerate}[label=\textbf{C\arabic*.}]
    \item \textbf{确认 automorphism：}  
    你的库里传入的 index / Galois element 是否真的实现了
    \[
    X^k \mapsto -X^k \, ?
    \]

    \item \textbf{确认 plaintext 构造：}  
    是否能构造
    \[
    \frac12,\qquad X^k,\qquad -X^{N-k}
    \]
    这三个明文乘子？

    \item \textbf{确认 level 对齐：}  
    \(\ctP\) 与 \(\ctm\) 在进入 \(\Add\) / \(\Sub\) 前是否处于同一 level？

    \item \textbf{确认 scale 对齐：}  
    \(\ctP\) 与 \(\ctm\) 的 scale 是否一致？  
    明文乘法后是否需要 rescale？

    \item \textbf{确认密文语义：}  
    你当前实现处理的是“槽语义 CKKS 流程”还是“本文采用的多项式语义流程”？不要把两者混用。
\end{enumerate}

\section{Conclusion}

本文第二版不再重复第一版的推导，而是将最终可用结论整理成了以下四个算法：

\begin{enumerate}[label=(\arabic*)]
    \item \textsc{ForwardSplitKeepY}：
    \[
    \ctP \longmapsto (\cte,\ctoy),
    \qquad
    \cte \rep \Pe(X^{2k}),\quad
    \ctoy \rep X^k\Po(X^{2k});
    \]

    \item \textsc{ForwardSplitNormalized}：
    \[
    \ctP \longmapsto (\cte,\cto),
    \qquad
    \cte \rep \Pe(X^{2k}),\quad
    \cto \rep \Po(X^{2k});
    \]

    \item \textsc{MergeFromKeepY}：
    \[
    (\cte,\ctoy)\longmapsto \ctP;
    \]

    \item \textsc{MergeFromNormalized}：
    \[
    (\cte,\cto)\longmapsto \ctP.
    \]
\end{enumerate}

就抽象结构而言，这四个算法已经是完整的；实现时只需把
\[
\Auto{1+2n},\ \Add,\ \Sub,\ \MulPt{\cdot}
\]
分别映射到具体库的 API，并处理好 automorphism key、plaintext 构造以及 level/scale 对齐即可。

\end{document}